% Complete binary helper compiler and matching base-two Pell proof.
\section{Binary coefficient codes and a matching Pell shift: 90 operations}
\label{sec:binary-product}

Changing the helper reverse rows makes every positive complementary
coefficient lie outside the indicator. Both fixed codes then have binary
digits. The geometric identity uses $\lambda$ directly in place of
$3\lambda$, saving one multiplication. A simultaneous change of the
main Pell base from $a+4$ to $a+2$ preserves all shared expressions.

\begin{theorem}\label{thm:best90}
Every recursively enumerable subset of the positive integers has an
admissible fixed index $(V,H,T_L)$ for a system with 34 positive unknowns
and 22 equations admitting a straight-line certificate of 90 operations:
48 multiplications and 42 additions. The supplied inputs are
$(x,V,H,T_L)$, with three fixed components besides the positive queried
input $x$. Fixed numerals and equality tests are free.
\end{theorem}

\subsection{The binary helper compiler}

Retain the physical representation of Section~\ref{sec:product-bound}.
The actual positive input $x$ has weight zero and is not split; the
homogeneous unit $\delta$ is one physical coordinate. Every other
logical value is a sum of three distinct nonnegative physical
coordinates, with disjoint groups for distinct logical values. The
fixed forbidden-bit decomposition of a value $z$ is $(z,0,0)$ if its
$H_0$ bit is zero, and $(z-H_0,H_0/2,H_0/2)$ otherwise.

Use the same two reserved helper triples $V_0,V_1$. Every ordinary
circuit row uses $V_0$ as input and mentions neither $x$ nor $V_1$.
Fresh square-difference copies handle repeated operands and repeated
unit factors. Expanded physical square coefficients are $\pm1$ and
cross coefficients are $\pm2$, so their divided coefficients are
$\pm1$.

For each $i=0,1$, retain the seed base row $-x^2$, augmented by its
own helper so that its tested value is $V_i^2-x^2$. Replace the
reverse row by
\begin{equation}\label{eq:binary-product-reverse}
 2x\delta-2V_i\delta.
\end{equation}
Its positive monomial has one non-input physical factor, while all
its negative monomials have two. Its augmenting helper is the other
$V_{1-i}$. The implication $V_i=x$ will be proved only after
$\delta>0$ has been recovered.

Include both signs of each ordinary zero equation, including
\begin{equation}\label{eq:binary-product-normalization}
\begin{gathered}
 X^2-V_0^2,\quad Y^2-V_0^2,\quad Z^2-V_0^2,\quad
 \delta'^2-\delta^2,\\
 2V_0X-2\delta\delta'-2u\delta.
\end{gathered}
\end{equation}
The logical groups are distinct. Keeping the extra copies $Y,Z$ is
convenient and has no effect on the final operation count. Finish with
two unpaired $\delta^2$ rows. Only the last is padded, by a preceding
raw coefficient $2xX$.

For $m$ true non-input physical coordinates, use
\[
 v_i=8\cdot7^i\quad(0\le i<m),\qquad M=8\cdot7^{m-1}.
\]
Sums of at most six such weights have unique multisets. If $s$ counts
the main rows, including the four helper rows and two unit rows, put
\begin{equation}\label{eq:binary-product-layout}
\begin{gathered}
 d_0=12M+8,\qquad
 t_j=(s+2)d_0+8M+jd_0\quad(0\le j<s),\\
 t_*=t_{s-1}=(2s+1)d_0+8M,\qquad K=t_*+3.
\end{gathered}
\end{equation}
Place each divided base coefficient $a$ of monomial weight $w$ at
$t_j-w$. For a negative coefficient at $r=t_j-w$, add the six
positive helper complements
\begin{equation}\label{eq:binary-product-complements}
 \sum_{h\le h',\ h,h'\in P}T^{r-v_h-v_{h'}},
\end{equation}
where $P$ is the own helper group for a seed, the other group for a
reverse row, and the $V_1$ group for an ordinary row. Let $\mathcal R$
contain every main target and every negative base exponent. The seed
target coincides with its negative exponent. Add one reset $T^{r-4}$
for each distinct $r\in\mathcal R$, and add the final padding
\begin{equation}\label{eq:binary-product-padding}
 D_{\rm pad}(T)=\sum_{h\in P(X)}T^{t_*-1-v_h}.
\end{equation}
Call the resulting polynomial $D(T)$.

The coefficient-isolation proof from Section~\ref{sec:product-bound}
has one new degree case. Helper complements in non-seed rows have
degree four, so cannot affect a main target after multiplication by
$C(T)^2$. At a negative position, another base monomial could
contribute only if
$w_{\rm base}=w_{\rm negative}+w_C$. Ordinary base monomials have
degree two; the new positive reverse monomial has degree one and
cannot satisfy this equality with a degree-two negative monomial.
At an extra negative-position target, an augmentation contributes
only if
\[
 w_{{\rm negative},1}+w_{\rm helper}
   =w_{{\rm negative},2}+w_C.
\]
The negative factors and helper group are disjoint. Multiset uniqueness
forces the same negative pair and helper pair, giving exactly
$V_{\rm helper}^2-x^2$. A seed separately gives $V_i^2-x^2$ at its
coincident target. Reverse main targets remain
\eqref{eq:binary-product-reverse}.

The same uniqueness and row spacing prevent coefficient accumulation.
Base and augmentation terms have residue zero modulo eight; resets
have residue four and padding has residue seven. All nonzero $D$
coefficients therefore remain $\pm1$. Define
\begin{equation}\label{eq:binary-product-codes}
 \ell_0(T)=\sum_{i=0}^{m-1}T^{v_i}
     +\sum_{r\in\mathcal R}(T^r+T^{r+1}+T^{r+2}),
 \qquad e_0(T)=\ell_0(T)+D(T).
\end{equation}

Every negative $D$ coefficient is at an indicator start. Every
positive coefficient is now outside the indicator. Ordinary positive
base complements are distinct degree-two monomials; reverse positive
complements have degree one; non-seed helper complements have degree
four; and seed helper complements differ from their sole main start.
Reset and padding residues are outside the three tested residues.
Other row bands and the low true-coordinate region are disjoint.
In particular, the previous positive $x^2$ term at the reverse main
target has disappeared. Thus both $\ell_0$ and $e_0$ have only
digits zero and one.

Both code polynomials have zero constant coefficient and support below
$K$. The leading $D$ term is the last reset $T^{t_*-4}$ with coefficient
$+1$; all other terms in the final positive-only row have smaller degree.
Consequently $D(B)>0$ for every integer $B\ge2$.

Permit dummy coordinate digits at the indicator windows. Their least
weight is at least $t_0-2M$, while the least $D$ exponent is at least
$t_0-4M$. Their sum exceeds $t_*+2$, so no dummy affects any tested
coefficient or reset. The least $D$ exponent also exceeds $M$, so
every true-coordinate indicator position has raw coefficient and
incoming carry zero.

\subsection{Fixed constants and bounds before Pell decoding}

Let $c_*$ count the input, true coordinates, and permitted dummies,
and let $D_1=\sum_h|[T^h]D(T)|$. Choose a power of two $L>3K+2$,
and then a power of two $H_0$ such that
\begin{equation}\label{eq:binary-product-threshold}
 H_0>\max\bigl(2\cdot2^{2L+1},
       128\max(1,D_1)c_*^2,\ 3L,\ 1024\bigr).
\end{equation}
The fixed index is
\begin{equation}\label{eq:binary-product-index}
 H=H_0-1,\qquad V=\ell_0(2)+2^Le_0(2),\qquad
 T_L=\psi_2(L).
\end{equation}
It is chosen before the queried input.

The source retains the positive product bound and uses the new alphabet:
\begin{equation}\label{eq:binary-product-source}
\begin{gathered}
 C=x+g,\qquad \sigma=(e-\ell)C^2,\qquad
 \ell+\sigma+\alpha=q,\\
 b=x+\beta,\qquad \theta=H+b=B-2,\qquad
 B=H_0+1+b,\qquad \lambda(B-1)=q^2-1.
\end{gathered}
\end{equation}
There is no supplied $\Omega$. Positivity gives $C\ge2$, $e-\ell\ge1$,
and therefore
\begin{equation}\label{eq:binary-product-ranges}
 \ell<q,\qquad C^2\le\sigma<q,\qquad
 e\le\ell+\sigma<q,\qquad g<C<q.
\end{equation}
Set
\begin{equation}\label{eq:binary-product-packing}
\begin{aligned}
 &S_2=\ell+eq=V+t\theta,\qquad n=q^8,\\
 &S=g+q^2(S_2+q^2\sigma),\\
 &T_+=q^2(1+\theta\lambda)+\ell(\theta q^4-b),\\
 &r=S(n^2-n)+T_+(n^2-1).
\end{aligned}
\end{equation}
Geometry gives $B\le q^2$, $\theta\lambda>b$, and
$1+\theta\lambda=q^2-\lambda<q^2$. Hence
\[
 T_+>bq^2-bq>0,\qquad
 T_+<q^4(1+\theta\ell)<Bq^5\le q^7,
 \qquad 0<S<q^7<n.
\]
Together with the fixed threshold, these imply
\begin{equation}\label{eq:binary-product-preliminary}
 n\ge64,\qquad n\le n^2-1\le r<2n^3,\qquad
 2\le L,\qquad 3L\le B\le n.
\end{equation}
No exponent or coefficient decoding has been used.

\subsection{The main Pell base is now \texorpdfstring{$a+2$}{a+2}}

Distinguish the Pell discriminant from the coefficient polynomial:
\begin{equation}\label{eq:binary-product-pell-definitions}
\begin{gathered}
 U=wn^2,\quad Y_P=s_Pn^2,\quad E=UY_P,\quad Q=UY_P^2,\quad
 P=2Q+1,\\
 a=Y_P(U+1),\qquad A=a+2,\qquad D_P=A^2-1,\qquad J=2r+1.
\end{gathered}
\end{equation}
Retain the first norm, positive interval, and index congruence:
\begin{equation}\label{eq:binary-product-first-pell}
\begin{gathered}
 \tau(\tau+1)=(E^2+U)(Y_Pk)^2,\\
 c=Y_Pk+\eta,\qquad k=\eta+\zeta,\qquad
 k=r+1+hE.
\end{gathered}
\end{equation}
The main norm, relaxed auxiliary norm, and half-parameter block are
\begin{equation}\label{eq:binary-product-auxiliary}
\begin{gathered}
 d^2=1+D_Pc^2,\qquad R=ic^2,\qquad
 K_P=R^2=D_P(f^2-1),\\
 u_*=J+jc=c+of,\qquad
 K_P(u_*^2-y_{\rm aux}^2)=1-y_{\rm aux}^2.
\end{gathered}
\end{equation}
Replace the first exponent congruence by
\begin{equation}\label{eq:binary-product-first-exponent}
 d=U+ac+\gamma(4a+3).
\end{equation}
The second norm, positive gap, fixed index, and exponent congruence are
\begin{equation}\label{eq:binary-product-second-pell}
\begin{gathered}
 c=\kappa+\phi,\qquad \mu^2=1+D_P\kappa^2,\qquad
 \kappa=T_L+\Delta a,\\
 \mu=q+\kappa(A-B)+\rho\bigl(D_P-(A-B)^2\bigr).
\end{gathered}
\end{equation}
Because $A=a+2$ and $B=\theta+2$, the shared difference is still
$A-B=a-\theta$.

\subsection{Exact indices before either exponential}

The preliminary bounds give
\[
 U,Y_P\ge n^2\ge4096,\qquad
 E\ge n^4>r+1,\qquad a>n^4>J.
\]
Classification of the first norm yields $k=\psi_P(t)$ at a positive
index. Its congruence modulo $E$ gives
$t=r+1+vE$ with $v\ge0$. Thus $t\ge r+1$. Since $P>A$ and
$c>Y_Pk>k$, comparison with the main norm
$c=\psi_A(p)$, $d=\chi_A(p)$ gives
\[
 p\ge t+1\ge r+2\ge66,\qquad c>AD_P^2,\qquad c>J.
\]
These are the generic hypotheses of the relaxed auxiliary and
half-parameter proofs above. They use only $A>1$, not its former
shift by four or evenness of $r$. Applying them to
\eqref{eq:binary-product-auxiliary} recovers
\[
 p=J,\qquad c=\psi_A(J),\qquad d=\chi_A(J).
\]

If $v\ge1$, ordinary Pell growth and $E>r$ give
\[
 \frac ck
 \le\frac{(2A)^{2r}}{(2P-1)^{r+E}}
 \le\left(\frac{2A}{2P-1}\right)^{2r}<\frac12,
\]
because the changed exact difference is
\begin{equation}\label{eq:binary-product-index-exclusion}
 (2P-1)-4A=4Y_P\bigl(U(Y_P-1)-1\bigr)-7>0.
\end{equation}
This contradicts $c/k>Y_P$. Consequently $k=\psi_P(r+1)$, still
without either exponential identity or a decoded circuit equation.

\subsection{Ratio growth, exponent recovery, and exact rounding}

Put $\xi=(U+1)^{2r}/U^r$. Standard lower and upper Pell bounds give
\begin{equation}\label{eq:binary-product-ratio}
\begin{aligned}
 \frac ck
 &\ge \xi\,
   \frac{(1+3/(2a))^{2r}}{(1+1/(2Q))^r}>\xi,\\
 \frac ck&<\xi(1+2/a)^{2r}.
\end{aligned}
\end{equation}
The strict lower bound uses
$(1+3/(2a))^2>1+3/a>1+1/(2Q)$, since $6Q>a$.
Also
$4r/a<8/n\le1/8<1/2$. The elementary estimate
$(1+t)^m\le(1-mt)^{-1}<1+2mt$ therefore gives
$c/k<\xi(1+8r/a)$.

The positive interval and the lower ratio bound imply $\xi<Y_P+1$.
Since $\xi>U^r$ and $Y_P$ is an integer,
\begin{equation}\label{eq:binary-product-ratio-growth}
 Y_P\ge U^r,\qquad a=Y_P(U+1)>U^{r+1}.
\end{equation}
Using $\xi<Y_P+1<2Y_P$ in the upper estimate gives
\begin{equation}\label{eq:binary-product-ratio-error}
 0<c/k-\xi<16r/(U+1).
\end{equation}
The right side is not yet asserted small.

The chi-congruence exponent criterion used above applies to
\eqref{eq:binary-product-first-exponent} with base two, since
$A-2=a$ and $4A-5=4a+3$. Its growth assumptions hold before its use:
\[
 2^{3J}=8\cdot64^r<U^{r+1}<a<A,\qquad
 U^3\le U^r<a<A.
\]
It follows that $U=2^J$. Since $U=wn^2$ and $n=q^8$, both $n$ and
$q$ are powers of two.

For the second index, the norm and gap in
\eqref{eq:binary-product-second-pell} give
$\kappa=\psi_A(t_0)$, $\mu=\chi_A(t_0)$ with $0<t_0<J$. Moreover,
$0<L<J$ follows from $3L\le B\le n\le r$. Reduction of the fixed
index congruence modulo $a$, using $A\equiv2\pmod a$, gives
$\psi_2(t_0)\equiv\psi_2(L)\pmod a$. Both quantities are strictly
smaller than $a$: for $0<t<J$,
\[
 0<\psi_2(t)\le4^{t-1}<2^{3J}<a.
\]
Thus they are equal as integers, and strict growth implies $t_0=L$.
The second exponent criterion has the required bounds
\[
 B^{3L}\le B^B\le n^n\le U^r<a<A,\qquad
 q^3\le n^n\le U^r<a<A.
\]
Its congruence in~\eqref{eq:binary-product-second-pell} therefore
gives $q=B^L$. As $q$ is a power of two, so is $B$.

Now $U=2^J=2\cdot4^r>32r$, so
\eqref{eq:binary-product-ratio-error} is less than $1/2$. Split the
binomial expansion as $\xi=F+\varepsilon$, where
\begin{equation}\label{eq:binary-product-binomial-split}
 F=\sum_{j=r}^{2r}\binom{2r}{j}U^{j-r},
 \qquad
 \varepsilon=\sum_{j=1}^{r}\frac{\binom{2r}{r-j}}{U^j}.
\end{equation}
Exact symmetry gives the sharper bound
\begin{equation}\label{eq:binary-product-tail}
 0<\varepsilon
 \le\frac{4^r-\binom{2r}{r}}{2U}<\frac14.
\end{equation}
Hence $F$ is the integer part of $\xi$ and
$F\equiv\binom{2r}{r}\pmod U$. The interval and ratio bounds yield
\[
 F<\xi<c/k<Y_P+1,\qquad
 Y_P<c/k<\xi+\tfrac12<F+\tfrac34.
\]
Integrality forces $Y_P=F$. Since $n^2$ divides $Y_P$ and $U$, it
divides $\binom{2r}{r}$. This establishes every hypothesis for the
packed binary masks without relying on a decoded coefficient row.

\subsection{Canonical codes, carry clearing, and the unit}

From $b<B<q$ and $\ell<q$, one has $b\ell<q^2$. Thus the exact
normalized mask is
\[
 T_+-1=(q^2-1-b\ell)+q^2\theta\lambda+q^4\theta\ell,
\]
with block widths $q^2,q^2,q^4$. The three independent masks follow
from central-binomial divisibility. The middle mask is
$(B-2)\sum_{j<2L}B^j$, so all digits of $S_2$ are binary. If $F_0(T)$
is its digit polynomial, the fixed congruence gives
$F_0(2)\equiv V\pmod{B-2}$. Both nonnegative quantities are below
$2^{2L}$, and~\eqref{eq:binary-product-threshold} makes their
difference smaller than $B-2$. Binary uniqueness therefore gives
\[
 F_0(T)=\ell_0(T)+T^Le_0(T),\qquad
 \ell=\ell_0(B),\qquad e=e_0(B).
\]
There is no ambiguity between the two codes because $\ell,e<q$.

The first mask has digit $H_0=B-1-b$ at each indicator position and
$B-1$ elsewhere. It confines $g$ to permitted digits and clears
their fixed bit. Its unit digit is zero; $x<B$ then makes $C=x+g$
exactly the physical polynomial evaluated at $B$. The third mask
requires binary digits of $\sigma=D(B)C(B)^2$ at every indicator.

Since $B\ge2H_0$, all raw coefficients satisfy
\[
 |[T^h]D(T)C(T)^2|
 \le D_1C(1)^2<D_1c_*^2B^2<B^3/256.
\]
The signed carry recurrence keeps incoming carries below $B^2/128$
in absolute value. Dummy exclusion makes the low raw residues
$0,4,7$ modulo eight. At a start $r$, empty positions $r-6,r-5$
reduce the carry entering the reset at $r-4$ to $-1$ or zero.
Its nonnegative raw polynomial contains $x^2\ge1$, so its outgoing
carry is nonnegative and below $B^2$. The empty positions $r-3,r-2$
reduce it to zero. Every $r-1$ is empty except the last unit pad,
and every $r+1,r+2$ is empty.

Every ordinary, helper, and extra target therefore tests its exact
raw value modulo $B^3$. Any negative nonzero value in the proved
range has top digit greater than one and fails. Use these facts in
the following order. The seeds give $V_i\ge x>0$. Both signs of
every ordinary row force it to vanish independently of the helper
values. In particular~\eqref{eq:binary-product-normalization} gives
\[
 X=Y=Z=V_0,\qquad \delta'=\delta,\qquad
 V_0^2=\delta^2+u\delta,
\]
so $\delta>0$. The reverse target $2\delta(x-V_i)$ cannot be
negative, proving $V_i\le x$ and hence $V_0=V_1=x$. All extra
targets now vanish, and the guard gives $0<\delta\le x<B$.

The final preceding coefficient from
\eqref{eq:binary-product-padding} is exactly $2xX=2x^2$. Its
incoming carry is zero. The two unit tests are consequently
\[
 \delta^2,\qquad \delta^2+\lfloor2x^2/B\rfloor
\]
with two empty following positions. The first gives
$\delta^2=kB+s$, $k,s\in\{0,1\}$. If $s=1$, then
$\delta^2\le B+1<B^2/16$, and the four roots of one modulo a
power of two leave only $\delta=1$ in $(0,B/4)$. Otherwise a
nonunit solution has $\delta^2=B$ and $x^2\ge B$.
Writing $h=\lfloor2x^2/B\rfloor$ gives $2\le h<2B$. Binary digits
of $B+h<3B<B^2$ would force
$1+\lfloor h/B\rfloor\le1$ and then $h\le1$, a contradiction.
Thus $\delta=1$, and the exact homogeneous circuit recovers the
represented system at its actual input.

\subsection{Positive necessity at the new Pell base}

Given a represented solution, choose $V_0=V_1=x$,
$\delta=\delta'=1$, $X=Y=Z=x$, $u=x^2-1$, all circuit values and
fresh copies, and their fixed forbidden-bit splits. Set all dummy
digits to zero. Every ordinary, seed, reverse, and extra target is
zero, and both unit targets are one.

Choose an arbitrarily large power of two $B$ after these physical
values have been fixed. Require $B>H_0+1+x$, every raw coefficient
in absolute value below $B/4$, and $2x^2<B/4$. Set
\begin{equation}\label{eq:binary-product-necessary-coding}
\begin{aligned}
 &b=B-H_0-1,\quad \beta=b-x,\quad g=C(B)-x,\\
 &\ell=\ell_0(B),\quad e=e_0(B),\quad q=B^L,\quad n=q^8,\\
 &\sigma=D(B)C(B)^2,\qquad \alpha=q-\ell-\sigma.
\end{aligned}
\end{equation}
The positive leading coefficient of $D$ and the unit coordinate give
$\sigma,g>0$. The degree bound $\deg(DC^2)<3K<L$ allows $B$ large
enough that $\alpha>0$. The first-mask bits clear, the middle digits
are binary, and every third window is $(0,0,0)$ or $(1,0,0)$; the
last pad has zero carry. Geometry supplies positive $\lambda$, and
the congruence quotient $(\ell+eq-V)/(B-2)$ is positive integral by
evaluation of the nonconstant nonnegative polynomial
$\ell_0(T)+T^Le_0(T)$ at $B$ and at two.

The masks give central-binomial divisibility for the resulting $r$,
with bounds~\eqref{eq:binary-product-preliminary}. Also $B$ divides
both $g$ and $\ell$, so~\eqref{eq:binary-product-packing} makes
$S,T_+,r$ even. Put $J=2r+1$ and construct the new Pell witnesses:
\begin{equation}\label{eq:binary-product-necessary-pell}
\begin{gathered}
 U=2^J,\quad w=U/n^2,\quad
 Y_P=\lfloor(U+1)^{2r}/U^r\rfloor,\quad s_P=Y_P/n^2,\\
 a=Y_P(U+1),\quad A=a+2,\quad E=UY_P,\quad Q=UY_P^2,\quad
 P=2Q+1,\\
 c=\psi_A(J),\quad d=\chi_A(J),\quad k=\psi_P(r+1).
\end{gathered}
\end{equation}
The quotient $w$ is positive integral because $n^2$ and $U$ are
powers of two and $n^2\le r^2<2^{2r+1}$. Binomial divisibility and
\eqref{eq:binary-product-tail} make $s_P$ positive integral.
The same ratio bounds at these chosen indices give
$Y_P<\xi<c/k<Y_P+3/4$. Thus
\[
 \eta=c-kY_P>0,\qquad \zeta=k-\eta>0.
\]
Set $\tau=(\chi_P(r+1)-1)/2$ and
$h=(k-r-1)/E$. The odd parameter $P$ makes $\tau$ integral;
congruence modulo $P-1=2Q$ gives integrality of $h$, and strict
Pell growth gives their positivity.

For the relaxed and half-parameter auxiliaries choose
\[
 m=2cJ,\quad f=\chi_A(m),\quad
 i=D_P\psi_A(m)/c^2,\quad R=ic^2,
\]
then take
\[
 y_{\rm aux}=\psi_R(J),\qquad
 u_*=\chi_R(J)/R,\qquad
 o=(u_*-c)/f,\qquad j=(u_*-J)/c.
\]
The generic positive construction of
Section~\ref{sec:half-parameter-pell} applies because even $r$
gives $J\equiv1\pmod4$. It proves every required divisibility and
strict positivity at the new value of $A$.

Finally choose $\kappa=\psi_A(L)$, $\mu=\chi_A(L)$,
$\phi=c-\kappa$, and
$\Delta=(\psi_A(L)-\psi_2(L))/a$. The inequalities $0<L<J$,
$A>2$, and $L\ge2$ give $\phi,\Delta>0$; polynomial congruence
modulo $a$ makes $\Delta$ integral. Both exponent congruences have
integer quotients. Their positivity follows explicitly from
\[
 d-ac=2c-\psi_A(J-1)>c>U,
\]
and
\[
 \mu-(A-B)\kappa
   =B\kappa-\psi_A(L-1)>(B-1)\kappa>q.
\]
Here $c,\kappa\ge A$, and the established growth gives
$A>U^3,q^3$. Both moduli are positive. This supplies every positive
unknown at the changed base; it does not retain old base-four
witnesses without justification.

\subsection{Exact operation count and verification}

Geometry now reads
\[
 q^2=(1+\theta\lambda)+\lambda.
\]
The former instruction $3\lambda$ is deleted. The shared exponent
modulus and discriminant are now $4a+3$ and $a^2+(4a+3)$ in place
of $8a+15$ and $a^2+(8a+15)$, with the same count; $A-B=a-\theta$
still uses one subtraction. Thus exactly one multiplication disappears:
48 multiplications and 42 additions.

The complete checker states all 22 source residuals afresh. If
$E_0=4a+12$ is the old discriminant minus the new one, their only
changes, in zero-based source order, are
\[
\begin{array}{c|l}
2&2\lambda\\
13&\gamma E_0\\
14&c^2E_0\\
15&(f^2-1)E_0\\
16&-(f^2-1)(u_*^2-y_{\rm aux}^2)E_0\\
17&\rho E_0\\
18&-\kappa^2E_0 .
\end{array}
\]
The three inherited residual corrections remain triangular, using
independently exact geometry and relaxed-norm equations. Every
primitive, equality, source difference, and supplied input is checked.
The generated table lists all 90 instructions and 22 free equalities.

The full compiler proof and exact certificate passed independent
review. The sparse regression checks a sample with 35 true coordinates,
24 main rows, 121 tested starts, 397 indicator digits, and 912
complementary coefficients. It verifies binary code digits, every
extra target, the exact pad, dummy exclusion, six canonical assignments,
four wide signed-carry assignments, and 1,252 local unit cases.
A separate exact Pell regression checks the changed ratio, norms,
congruences, and positive witnesses through indices $r=64,66$.
These finite receipts supplement the general proofs and do not assert
exhaustive coverage or optimality.
